% Generated draft: requires mathematical and source review.
\documentclass[11pt]{article}
\usepackage[margin=25mm]{geometry}
\usepackage{fontspec}
\setmainfont{Latin Modern Roman}
\usepackage{amsmath,amssymb,mathtools,unicode-math}
\setmathfont{Latin Modern Math}
\usepackage{enumitem,needspace,xurl,hyperref}
\setlength{\parindent}{0pt}
\setlength{\parskip}{5pt}
\setlength{\emergencystretch}{3em}
\newcommand{\N}{\mathbb{N}}
\newcommand{\Z}{\mathbb{Z}}
\newcommand{\Q}{\mathbb{Q}}
\newcommand{\R}{\mathbb{R}}
\newcommand{\C}{\mathbb{C}}
\newcommand{\F}{\mathbb{F}}
\newcommand{\A}{\mathbb{A}}
\newcommand{\PP}{\mathbb{P}}
\newcommand{\E}{\mathbb{E}}
\newcommand{\eps}{\varepsilon}
\newcommand{\dd}{\,\mathrm{d}}
\begin{document}
\section*{UnsolvedMath ID: 801070}
\textbf{Source problem code:} catalogue entry
\subsection{Definitions and mathematical statement}
The supplied Question 22 asks two questions under the hypotheses that \(f\) is quasi-hyperbolic and dissipative and that \(J\) is connected. The excerpt does not define the ambient class of maps, \(J\), quasi-hyperbolicity, external rays, or a finite quotient of the real solenoid. Consequently, the following mathematical interpretation is conditional and requires comparison with the cited source.

In the customary complex H\'enon setting, one considers a polynomial automorphism \(f:\C^2\to\C^2\) of H\'enon type, of dynamical degree \(d\geq 2\). In this setting, define
\[
K^{\pm}(f)=\left\{z\in\C^2:\{f^{\pm n}(z):n\geq 0\}\text{ is bounded}\right\},\qquad
J^{\pm}(f)=\partial K^{\pm}(f),\qquad
J(f)=J^+(f)\cap J^-(f).
\]
Dissipativity means \(0<|\det Df|<1\), where the Jacobian determinant is constant. These conventions are not specified in the supplied excerpt. The hypothesis of quasi-hyperbolicity must retain its meaning in the source; no definition of that property is supplied here.

Under the source's intended conventions, ask whether every such \(f\) with connected \(J(f)\) has each of the following properties:
\begin{enumerate}
\item Every external ray lands at a point of \(J(f)\). More explicitly, if an external ray is parametrized as \(\gamma:(0,\infty)\to\C^2\), with \(t\downarrow 0\) corresponding to its end toward \(J(f)\), the requested landing property is
\[
\exists p_\gamma\in J(f)\quad \lim_{t\downarrow 0}\gamma(t)=p_\gamma.
\]
This formula describes landing; it does not define which curves are the external rays in the source.
\item The space \(J(f)\) is a finite quotient of the real solenoid, in the sense intended by the source.
\end{enumerate}
For clarity about the second unresolved terminology, a common degree-\(d\) solenoid is the inverse limit
\[
\mathcal S_d=\left\{(\zeta_0,\zeta_1,\ldots)\in(S^1)^{\N_0}:\zeta_{n+1}^{\,d}=\zeta_n\text{ for every }n\geq 0\right\},\qquad
S^1=\{\zeta\in\C:|\zeta|=1\},
\]
with the subspace topology from the product, where \(\N_0=\{0,1,2,\ldots\}\). If the intended solenoid is \(\mathcal S_d\) and a finite quotient means a finite-to-one continuous image, the second question asks for a continuous surjection
\[
q:\mathcal S_d\longrightarrow J(f),\qquad
\#q^{-1}(p)<\infty\quad\text{for every }p\in J(f).
\]
The excerpt does not establish this interpretation, identify the solenoid's degree, or specify whether uniformly bounded fibers, a group quotient, or compatibility with the dynamics is required. These qualifications remain unresolved.
\subsection{Short English statement}
Question 22 asks whether quasi-hyperbolicity, dissipativity, and connectedness of \(J\) force every external ray to land on \(J\), and whether \(J\) is a finite quotient of the real solenoid. The ambient setting and several essential definitions are absent from the excerpt; the displayed complex H\'enon interpretation is provisional.
\subsection{Sources}
\begin{itemize}
\item Local UnsolvedMath registry, supplied complete record: permanent ID \(801070\), source problem code \texttt{catalogue entry}, title \emph{H\'enon Maps: Open Problems --- Question 22}. The registry attributes the problem to Berger and Bedford, with proposed year 2023, and labels it open. This label is supplied curation, not an independently verified status assessment.
\item Berger and Bedford, \emph{H\'enon Maps: A List of Open Problems} (2023), Question 22. Supplied source URL: \url{https://arxiv.org/abs/2312.03907}. The background records access on 2026-10-16 and extraction by \texttt{arxiv\_latex}; these are supplied provenance statements, not access or verification performed for this draft.
\end{itemize}
\end{document}
